% SPDX-License-Identifier: Apache-2.0
% covers: EthRx
\datasheet{EthRx}{The receive half of the Ethernet MAC: frames off GMII, checked, onto a channel.}

\dsfacts{\code{lib/parts/src/eth.rs}}{\code{txhdl\_parts::eth::EthRx}}%
{\code{//docs:eth}}

\subsection*{Function}

\code{EthRx} takes a frame off GMII: \code{rxd}, a byte per cycle,
\code{rx\_dv} and \code{rx\_er}. It stores the whole frame and checks
its CRC-32 frame check sequence. Only then does it offer the frame's
bytes on the channel \code{rx}, each an \code{EthByte}, without the
preamble and the check sequence. It holds two frames, so it takes the
next frame off the wire while it still offers the last (issue 1409).
It drops a frame whose check fails, and counts every frame it drops
by cause; the counts leave on \code{rx\_counts} in Gray code, for
another clock to read (issue 1404).
While it offers a frame, \code{rx\_len} says how many bytes the frame
has, so a consumer that must know the length before it has taken the
frame can read it; \code{EthSlots} does.
The unit runs on one clock, so on a board it can run on the receive
clock the PHY recovers from the line.

\subsection*{Parameters}

\code{EthRx} has no type parameters. Its store holds
\code{RX\_STORE} bytes, 4096: two frames of \code{FRAME\_MAX}, 2048,
the second from byte 2048. The store says
\code{ram\_style("block")}, and is read only into the register
\code{rq}, so it is one block RAM (issue 1421).

\dstables{EthRx}

\subsection*{Behaviour}

The unit is two processes. The receiver watches the wire every cycle;
the offerer is the sequence of giving a frame away: wait for a full
frame, then its bytes without the check sequence, one per turn, each
held until the consumer takes it. The lowering numbers the offerer's
waits into a state register the unit does not declare, \code{at\_wait},
keeps its loop's count in \code{for0}, and holds \code{rx\_len} between
states in \code{rx\_len\_held}; all are in the tables above.

\begin{itemize}
\item The receiver stores into the frame \code{wsel} names, and the
offerer offers the frame \code{rsel} names; \code{full0} and
\code{full1} say which hold a good frame, and \code{len0} and
\code{len1} their lengths. \code{olen} is the length of the frame
whose turn it is to be offered, which the receiver keeps.
\item While hunting, that is while it is not receiving, the frame
\code{wsel} names is not full, and it is not skipping, the receiver
starts a frame on \code{rx\_dv} high with \code{0xd5} on \code{rxd}:
it sets \code{receiving}, clears \code{len}, \code{bad} and
\code{long}, and sets the CRC register to \code{CRC\_INIT}. It passes
over \code{0x55}.
\item Any other byte with \code{rx\_dv} high while hunting is the middle
of a frame whose start the unit missed. The receiver sets \code{skip}
and ignores the line until \code{rx\_dv} falls.
\item While receiving, the receiver stores a byte per cycle for as long
as \code{rx\_dv} stays high, into the frame \code{wsel} names. The
CRC register runs over every byte, the check sequence included.
\code{bad} records any cycle with \code{rx\_er} high, and \code{long}
a byte past the frame's 2048th.
\item When \code{rx\_dv} falls, the frame is good if \code{bad} and
\code{long} are low, the register holds \code{CRC\_RESIDUE}, and
\code{len} is more than four. A good frame sets its full bit and its
length, and \code{wsel} turns to the other frame. The receiver drops
any other frame and adds one to \code{dropped}, and to
\code{drop\_size} if it was four bytes or fewer or long, or else to
\code{drop\_check}.
\item The full frame \code{rsel} names starts the offerer.
\code{rx\_len} reads the frame's length without the check sequence
while it is offered, and zero at every other time. The offerer sets
\code{ri} one before the frame's first byte; each step then moves
\code{ri} on and reads the byte there into \code{rq}, and offers
\code{rq} on \code{rx} until \code{rx} takes it, marking the last of
\code{olen} less four. The step comes at the edge a byte is taken, so
a frame leaves a byte a cycle to a consumer that takes one each cycle.
Then \code{given} is up for one cycle, which adds one to
\code{frames}, clears that frame's full bit, and turns \code{rsel} to
the other frame.
\item A frame that arrives while both frames are full is dropped: the
receiver sets \code{skip} and adds one to \code{dropped} and to
\code{drop\_room}, once for that frame.
\item \code{counts} holds \code{drop\_room}, \code{drop\_check} and
\code{drop\_size} in Gray code, sixteen bits each, room highest, as
\code{gray16} writes them, and \code{rx\_counts} reads it. A register,
so each count changes in one bit at an edge, and a clock unrelated to
this one can sample it through two flip-flops; \code{ungray16} decodes
it there.
\end{itemize}

\subsection*{Verification}

\begin{itemize}
\item The tests below are in \code{lib/parts/src/eth.rs}, and run in
\code{//lib/parts:parts\_test}. Each loops \code{EthTx} into the
receiver a cycle late.
\item \code{a\_frame\_goes\_out\_as\_the\_model\_says\_and\_comes\_back}
checks that a frame of 20 bytes comes back padded to sixty, and a frame
of 100 bytes comes back as it went.
\item \code{frames\_at\_the\_minimum\_length} does the same at 59, 60
and 61 bytes.
\item \code{a\_corrupted\_frame\_is\_dropped} flips a byte of the second
of three frames. It checks that the other two come back and that one
frame is counted dropped.
\item \code{a\_frame\_on\_a\_busy\_receiver\_is\_dropped} holds the
client off. It checks that the third of three frames is dropped and
counted, and that the first two are intact.
\item The tests below drive the receiver's wire directly, as a sender
puts frames on it, with twelve idle bytes between them, and read at
the board's pace, four bytes in five cycles, unless they say
otherwise.
\item \code{two\_full\_frames\_at\_the\_minimum\_gap\_both\_come\_through}
sends two frames of 1514 bytes and checks that both come out and none
is dropped (issue 1409).
\item \code{a\_train\_at\_the\_minimum\_gap\_loses\_nothing\_uncounted}
sends six, and checks that every frame comes out whole and in order or
is counted.
\item \code{the\_drops\_are\_counted\_by\_cause} breaks frames on the wire
and holds the reader off, and checks the three counts both in their
registers and decoded from \code{rx\_counts} (issue 1404).
\item \code{a\_count\_survives\_gray\_code} takes every sixteen-bit
count through \code{gray16} and \code{ungray16} and checks that each
step changes one bit.
\item \code{the\_receivers\_store\_is\_read\_as\_a\_block\_ram} checks
that the netlist marks the store a block RAM, reads it into \code{rq},
and reads none of it straight onto \code{rx}; the lowering refuses such
a memory reached at more ports than a block RAM has (issue 1421).
\item \code{ex\_eth} runs the unit with one byte of the second frame
flipped on the wire, and asserts that the receiver dropped one frame,
for its check.
The build simulates the netlist of \code{EthRx} against that run under
nvc and Verilator: \code{//docs:eth\_sim\_eth\_rx\_eth\_rx\_tb\_test}
and \code{//docs:eth\_vsim\_eth\_rx\_test}.
\end{itemize}

\subsection*{Used by}

The example \code{ex\_eth}, and \code{ex\_ethdma} and \code{ex\_remote}
ahead of their receiving sides. The flagship: \code{//eth:mac\_v} writes
the Verilog of \code{eth\_rx}, and \code{flagship/board/flagship.v}
instantiates it on the PHY's receive clock, crosses its bytes to the
core's clock through \code{chan\_cdc}, and samples \code{rx\_counts}
there through two flip-flops into the board's \code{mac\_drops}, which
\code{EthSlots} reads at words 14 and 15. The echo design for the
AX7A200B:
\code{//eth:netlist} writes the Verilog of \code{eth\_rx}, and
\code{eth/board/eth\_echo.v} instantiates it on the PHY's receive
clock. \code{eth/board/eth\_cdc.v} takes its bytes across to the
transmit clock. \code{//eth:echo\_synth} and \code{//eth:echo\_pnr} put
that design through Vivado to a bitstream. \code{//docs:eth} states
that it has not yet run on the board.

\subsection*{Limits}

The unit holds two frames. A client that has taken neither loses the
next one, which is counted. A consumer slower than the wire cannot keep
up with an unbroken train of frames, since two frames are all the room
there is. \code{len} stops at 2047, so the unit stores any further
bytes of a longer frame at the frame's last byte, and drops the frame
for its size. The counts are sixteen bits and wrap. Only gigabit works:
at 100 or 10~Mbit a PHY sends a nibble per cycle, and the unit does not
take a byte over two cycles. The unit reads \code{rx\_er} only while
receiving.
